\section{Éléments de preuve complémentaires en annexe}

Les documents de cadrage détaillés complètent l'analyse du besoin et la définition du périmètre présentées au début de ce chapitre. Le cahier des charges complet est reproduit en annexe~\ref{ann:ipf-cahier-charges}. Les personae et leurs parcours sont détaillés en annexe~\ref{annexe:persona}, tandis que les exigences fonctionnelles et non fonctionnelles sont respectivement accessibles en annexes~\ref{ann:ipf-exigences-fonctionnelles} et~\ref{ann:ipf-exigences-non-fonctionnelles}. La matrice de synthèse associée est reproduite en annexe~\ref{ann:ipf-matrice-exigences}.

La trajectoire temporelle décrite dans la section consacrée à la feuille de route peut être examinée dans son intégralité grâce au calendrier directeur en annexe~\ref{ann:ipf-calendrier-directeur} et au diagramme de Gantt complet en annexe~\ref{ann:ipf-gantt-complet}. Ces pièces permettent de retrouver les phases, les jalons et les dépendances sans surcharger le corps du dossier.

Le dispositif de suivi est documenté plus précisément par l'annexe consacrée aux indicateurs de projet (annexe~\ref{annexe:kpi-projet}). Elle distingue les indicateurs effectivement exploitables pendant le projet de ceux qui nécessitent une recette ou une exploitation réelle. Le registre exhaustif des risques, avec probabilités, impacts, scores et plans de réponse, est fourni en annexe~\ref{annexe:matrice-risque}.

La démarche de validation fonctionnelle est détaillée dans les fiches de tests reproduites en annexe~\ref{ann:ipf-tests-fonctionnels}. Elles complètent le bilan du produit en explicitant les préconditions, scénarios, résultats attendus et observations pour les parcours apprenant et administrateur.

Enfin, la trame utilisée pour les entretiens individuels de fin de projet est conservée en annexe~\ref{annexe:entretiens-fin-projet}. Elle complète l'analyse du fonctionnement collectif en montrant les thèmes abordés : ressenti, rôle, compétences mobilisées, blocages, coordination, apprentissages et pistes d'amélioration.